\documentclass[11pt]{article}

\usepackage[margin=1.02in]{geometry}
\usepackage{amsmath,amssymb,amsthm,mathtools}
\usepackage{booktabs,tabularx,array}
\usepackage{enumitem}
\usepackage{microtype}
\usepackage{xcolor}
\usepackage[hidelinks]{hyperref}
\usepackage{fancyhdr}

\definecolor{deepblue}{HTML}{173B57}
\hypersetup{
  colorlinks=true,
  linkcolor=deepblue,
  citecolor=deepblue,
  urlcolor=deepblue,
  pdftitle={Exact Quantum Values and Permutation-Blind Maximizers in Cyclic Bell Inequalities},
  pdfauthor={Alec Kriebel},
  pdfsubject={Cyclic Bell inequalities, exact quantum values, and global randomness},
  pdfkeywords={Bell inequalities, Tsirelson bounds, commuting operators, device-independent randomness}
}

\pagestyle{fancy}
\fancyhf{}
\fancyhead[L]{\small\textsc{Cyclic Bell values and maximizers}}
\fancyhead[R]{\small\textsc{A. Kriebel}}
\fancyfoot[C]{\thepage}
\setlength{\headheight}{14pt}

\newtheorem{theorem}{Theorem}[section]
\newtheorem{lemma}[theorem]{Lemma}
\newtheorem{proposition}[theorem]{Proposition}
\newtheorem{corollary}[theorem]{Corollary}
\theoremstyle{definition}
\newtheorem{definition}[theorem]{Definition}
\theoremstyle{remark}
\newtheorem{remark}[theorem]{Remark}

\newcommand{\Id}{\mathbf 1}
\newcommand{\cI}{\mathcal I}
\newcommand{\cF}{\mathcal F}
\newcommand{\cJ}{\mathcal J}
\newcommand{\cH}{\mathcal H}
\newcommand{\Reop}{\operatorname{Re}}
\newcommand{\Tr}{\operatorname{Tr}}
\newcommand{\supp}{\operatorname{supp}}
\newcommand{\spec}{\operatorname{spec}}
\newcommand{\e}{\mathrm e}
\newcommand{\bbC}{\mathbb C}
\newcommand{\bbT}{\mathbb T}
\newcommand{\bbZ}{\mathbb Z}
\newcommand{\absop}[1]{\lvert #1\rvert}

\title{\vspace{-1.35cm}
\textbf{Exact Quantum Values and Permutation-Blind Maximizers}\\[-0.05em]
\textbf{in Cyclic Bell Inequalities}\\[0.45em]
\large Sharp operator bounds, equality structure, and limits of Bell-value randomness certification}
\author{Alec Kriebel\\
\small Independent Researcher}
\date{8 August 2026}

\begin{document}
\maketitle

\begin{abstract}
Perito, D'Avino, Jung, Mironowicz, Ac\'in, and Augusiak introduced two
cyclic Bell families designed for efficient certification of maximal global
randomness and conjectured the exact value of the reduced operator underlying
the first family.  We prove the conjectured value in every dimension and
strengthen the upper bound to the commuting-operator model.  The proof uses a
polar-decomposition positive factor and a sharp scalar extremum on the unit
circle.  We then study the equality mechanism and identify a permutation
freedom invisible to every first-harmonic correlator in the Bell score.  For
every \(d\geq4\), a final-two phase swap gives admissible order-\(d\),
rank-one projective measurements that attain the exact maxima of both
augmented cyclic families while producing a nonuniform designated joint
output distribution.  Hence the maximum scalar Bell value alone does not
certify \(2\log_2d\) bits of global randomness across the whole maximizing
face.  This does not challenge randomness certification conditioned on the
complete canonical behavior, nor does it show that the canonical strategy
lacks maximal randomness.  We give an exact \(d=4\) certificate over
\(\mathbb Q(\zeta_{16})\), all-dimensional analytic constructions, and
scoped consequences for low-setting designs.
\end{abstract}

\medskip
\noindent\textbf{Status.}
This is an unrefereed, AI-assisted manuscript.  Its proofs and verification
artifacts are public for specialist audit; internal adversarial review and
passing code are not peer review.

\section{Introduction}

Bell inequalities can certify randomness without a model of the devices, but
the conditioning data matter.  A complete behavior
\(p(a,b\mid x,y)\) contains far more information than one scalar Bell value.
If a Bell maximum exposes a unique behavior, symmetries of the score may force
uniform outputs and, with an operator-level argument, privacy against a
purifying adversary.  If the maximizing face contains several behaviors, the
same symmetry argument need not apply to each of them
\cite{DharaEtAl2013,NietoSillerasEtAl2014,BancalEtAl2014}.

Perito, D'Avino, Jung, Mironowicz, Acin, and Augusiak recently introduced two
families built around cyclic Fourier correlators
\cite{PeritoEtAl2026}.  They defined the families, found explicit
maximally-entangled strategies, proved the exact value of the second reduced
family by a sum-of-squares identity, and conjectured the exact value of the
first reduced family.  Their symmetry analysis motivated maximal-global-
randomness statements; for selected dimensions they also evaluated Eve's
guessing probability numerically after fixing the complete canonical
behavior.  These are logically different inputs to a randomness problem.

The surrounding literature includes high-dimensional and modular Bell
inequalities \cite{BuhrmanMassar2005,SalavrakosEtAl2017,KaniewskiEtAl2019},
self-testing from few measurements \cite{SarkarEtAl2021,KlepEtAl2026}, and
analytic or algorithmic SOS constructions
\cite{BarizienEtAl2024,CocciaEtAl2026}.  Two contemporaneous works of the
originating collaboration pursue all-dimensional global randomness and
three-outcome noise robustness by different constructions
\cite{FarkasMironowiczAugusiak2026,DavinoEtAl2026}.  We use these works to
set the attribution boundary; the family-specific polar bound and the
phase-permuted maximizers below are not imported from them.
To our knowledge, no prior work establishes either the first-family upper
bound in this normalization or the nonuniform exact maximizers constructed
here; the accompanying priority audit records the qualified search basis for
that statement.

This paper unifies the exact-value and randomness questions through the
equality phases of the first operator.  The polar proof reduces its upper
bound to a scalar roots-of-unity extremum.  At equality the relative unitary
may carry the full set of maximizing phases in different cyclic orders.
Those orders leave the Bell-visible first harmonics unchanged, but they alter
higher Fourier data and can change a complete output table.  The same Bob
observables also annihilate the published sum-of-squares factors for the
second family.

\begin{table}[t]
\centering
\small
\begin{tabularx}{\textwidth}{@{}>{\raggedright\arraybackslash}p{0.22\textwidth}
>{\raggedright\arraybackslash}X>{\raggedright\arraybackslash}X@{}}
\toprule
Topic & Originating work \cite{PeritoEtAl2026} & This paper\\
\midrule
Bell families & Introduced both cyclic families and their normalization &
Adopts the displayed Hermitian operator definitions, with a convention note\\
Attaining strategies & Supplied the canonical strategies and matching lower
bounds & Rechecks admissibility and attainment exactly\\
First reduced value & Conjectured \(2\csc(\pi/(2d))\); numerical evidence
through \(d=6\) & Proves the all-dimensional upper bound and the
\(q=qa=qc\) equality\\
Second reduced value & Proved value \(d\) by an SOS identity & Uses that SOS,
including its commuting-operator reading\\
Equality phases & Not a classification of the maximizing face & Gives a
conditional phase-permutation mechanism and its exact product hypotheses\\
Alternative maximizers & Not identified & Constructs biased exact maximizers
of both augmented families for every \(d\geq4\)\\
Canonical full behavior & Randomness studied numerically in selected
dimensions; second family self-tested at \(d=3\) & Not contradicted\\
Scalar-value implication & Maximal-violation certification proposed, with a
uniqueness-based symmetry rationale & Shows that the scalar maximum alone is
insufficient for \(d\geq4\)\\
\bottomrule
\end{tabularx}
\caption{Attribution and logical scope.  The present work does not claim the
family definitions, canonical strategies, or second-family SOS.}
\label{tab:contributions}
\end{table}

Our main conclusions are as follows.
\begin{enumerate}[leftmargin=1.6em]
\item The finite-dimensional tensor-product, approximate, and
commuting-operator values of the first reduced cyclic operator are all
\(2\csc(\pi/(2d))\).
\item A polar-linear permutation theorem gives explicit sufficient conditions
under which phase permutations remain admissible maximizers and preserve the
entire first-harmonic correlator matrix.
\item For each \(d\geq4\), one such maximizer has uniform local marginals but
a nonuniform designated joint table.  With trivial Eve this already gives a
guessing probability strictly larger than \(1/d^2\).
\item The same construction maximizes the second augmented cyclic family and
has the same target table.
\item The result separates scalar Bell-value certification from certification
conditioned on a fixed full behavior.  It is not a claim about the exact
worst-case guessing probability, the complete maximizing face, or the
randomness of the canonical realization.
\end{enumerate}

The equality mechanism is proved in Sections~\ref{sec:exact}--\ref{sec:phase};
the two biased families are treated in Sections~\ref{sec:biased} and
\ref{sec:second}; and the randomness interpretation is isolated in
Section~\ref{sec:randomness}.  Section~\ref{sec:settings} records only the
setting-complexity consequences that survive an independent audit.

\section{Bell operators, behaviors, and randomness notions}
\label{sec:framework}

Fix \(d\geq2\), write \(\omega=\e^{2\pi i/d}\), and identify outcomes with
\(\bbZ_d\).  A projective \(d\)-outcome measurement
\(\{M_a\}_{a\in\bbZ_d}\) is encoded by the order-\(d\) unitary
\[
 A=\sum_{a=0}^{d-1}\omega^aM_a,
 \qquad
 M_a=\frac1d\sum_{k=0}^{d-1}\omega^{-ak}A^k.
\]
Zero spectral projections are allowed in the general definition.  All
explicit strategies below have the complete simple spectrum of \(d\)-th
roots, so their outcomes are rank-one projectors.

For finite-dimensional tensor-product strategies, let \(q\) denote the
supremum of a Bell functional, and let \(qa\) denote the value over the
closure of those correlations.  In the commuting-operator model \(qc\),
Alice's and Bob's PVMs act on one Hilbert space and commute across parties.
We use
\[
 \mathcal Q_q\subseteq\mathcal Q_{qa}\subseteq\mathcal Q_{qc}.
\]
The operator inequalities below hold in the commuting model; a matching
finite-dimensional strategy then identifies all three values.

For a behavior \(p(a,b\mid x,y)\), a scalar Bell maximum fixes only one
supporting hyperplane of the quantum correlation set.  We distinguish:
\begin{itemize}[leftmargin=1.5em]
\item the value of a Bell functional;
\item its full maximizing face;
\item a selected canonical strategy and its canonical behavior;
\item the complete observed behavior, including higher Fourier correlators;
\item uniform observed target probabilities; and
\item private randomness against a purifying adversary.
\end{itemize}

For a purification \(\lvert\Psi\rangle_{ABE}\) and target settings
\((x_*,y_*)\), let
\[
 \sigma_E^{ab}
 =\Tr_{AB}\!\left[(M_a^{x_*}\otimes N_b^{y_*}\otimes\Id)
 \lvert\Psi\rangle\!\langle\Psi\rvert
 (M_a^{x_*}\otimes N_b^{y_*}\otimes\Id)\right].
\]
Eve's guessing probability is
\[
 G(AB\mid x_*,y_*,E)
 =\max_{\{Q_{ab}\}}\sum_{a,b}\Tr(Q_{ab}\sigma_E^{ab}).
\]
Maximal private global randomness means
\(\sigma_E^{ab}=\rho_E/d^2\) for every \(a,b\), and implies
\(G=1/d^2\).  Uniform probabilities are necessary but not sufficient for
this operator-valued condition.  Conversely, one explicit nonuniform target
table, even with one-dimensional Eve, is enough to rule out maximal private
randomness for that realization.

\begin{remark}[Normalization convention]
\label{rem:normalization}
The current source version of Ref.~\cite{PeritoEtAl2026} has a localized
normalization discrepancy: the displayed reduced value in its Eq.~(17) and
its \(d=3\) augmented value agree with the operator definition, while the
formula printed beside its first augmented randomness conjecture contains an
additional factor \(d\) in the denominator.  We take the displayed Hermitian
operator as authoritative.  Thus the reduced and augmented values in our
convention are respectively
\[
 M_d=2\csc\!\left(\frac{\pi}{2d}\right),
 \qquad M_d+1.
\]
Appendix~\ref{app:conventions} records the corresponding adjoint conventions.
Nothing in our argument uses the typographical discrepancy.
\end{remark}

\section{Exact value of the first cyclic operator}
\label{sec:exact}

For commuting Alice and Bob unitaries define
\begin{equation}
 \cI_d
 =\sum_{y=0}^{d-1}
 \Reop\!\left[(A_0+\omega^yA_1)B_y\right].
 \label{eq:Id}
\end{equation}
In the tensor model the product in \eqref{eq:Id} means a tensor product.  The
order-\(d\) relations are part of the Bell scenario, but will not be needed
for the upper bound.

\begin{theorem}[Sharp commuting-operator value]
\label{thm:exact}
Let \(A_0,A_1,B_0,\ldots,B_{d-1}\) be unitaries on an arbitrary complex
Hilbert space and suppose \([A_x,B_y]=0\) for every \(x,y\).  Then
\begin{equation}
 \cI_d\leq M_d\Id,
 \qquad M_d=2\csc\!\left(\frac{\pi}{2d}\right).
 \label{eq:main-bound}
\end{equation}
After imposing the order-\(d\) measurement relations,
\[
 \boxed{\beta_q(\cI_d)=\beta_{qa}(\cI_d)=\beta_{qc}(\cI_d)=M_d}
 \qquad(d\geq2).
\]
\end{theorem}

The proof has two independent ingredients.  The first is an operator identity
that remains valid when a polar factor has a kernel.

\begin{lemma}[Polar positive-factor identity]
\label{lem:polar}
Let \(\mathcal A,\mathcal B\subseteq\mathcal B(\cH)\) be commuting unital
\(*\)-algebras, let \(C\in\mathcal A\), let
\(C=V\absop C\) be its canonical polar decomposition, and let
\(B\in\mathcal B\) be unitary.  Put
\[
 P_{C,B}=\absop{C^\dagger}^{1/2}-V\absop C^{1/2}B.
\]
Then
\begin{equation}
 \frac{\absop C+\absop{C^\dagger}}2-\Reop(CB)
 =\frac12P_{C,B}^\dagger P_{C,B}.
 \label{eq:polar-factor}
\end{equation}
\end{lemma}

\begin{proof}
The initial projection \(V^\dagger V\) is the support projection of
\(\absop C\), while polar functional calculus gives
\[
 \absop{C^\dagger}^{1/2}V=V\absop C^{1/2},
 \qquad
 \absop{C^\dagger}^{1/2}V\absop C^{1/2}=C.
\]
The polar operators lie in the von Neumann algebra generated by
\(\mathcal A\), and hence commute with \(B\).  Expanding yields
\[
 P_{C,B}^\dagger P_{C,B}
 =\absop{C^\dagger}+\absop C-CB-B^\dagger C^\dagger.
\]
No inverse and no unitary extension of \(V\) has been used.
\end{proof}

The second ingredient is a scalar equality calculation; its elementary
parity-split proof is included in Appendix~\ref{app:scalar}.

\begin{lemma}[Roots-of-unity extremum]
\label{lem:scalar}
For \(z\in\bbT\),
\begin{equation}
 \sum_{y=0}^{d-1}\absop{1+\omega^yz}\leq M_d.
 \label{eq:scalar-bound}
\end{equation}
Equality holds exactly when
\begin{equation}
 z^d=(-1)^{d-1}.
 \label{eq:scalar-equality}
\end{equation}
\end{lemma}

\begin{proof}[Proof of Theorem~\ref{thm:exact}]
Set
\[
 L_y=A_0+\omega^yA_1,
 \qquad U=A_0^\dagger A_1,
 \qquad M_y=\Id+\omega^yU.
\]
Then \(L_y=A_0M_y\).  The operator \(M_y\) is normal because
\[
 M_y^\dagger M_y=M_yM_y^\dagger
 =2\Id+\omega^yU+\omega^{-y}U^\dagger.
\]
Consequently
\[
 \absop{L_y}=\absop{M_y},
 \qquad
 \absop{L_y^\dagger}=A_0\absop{M_y}A_0^\dagger.
\]
Put
\[
 F_d(U)=\sum_{y=0}^{d-1}\absop{\Id+\omega^yU}.
\]
Continuous functional calculus for the unitary \(U\), applied to
Lemma~\ref{lem:scalar}, gives \(0\leq F_d(U)\leq M_d\Id\).  Summing
Lemma~\ref{lem:polar} now gives
\[
 \cI_d
 \leq\frac12\bigl(F_d(U)+A_0F_d(U)A_0^\dagger\bigr)
 \leq M_d\Id.
\]
This uses only cross-party commutation, not a tensor-product factorization.
The finite order-\(d\) strategy rechecked in
Appendix~\ref{app:attainment} attains \(M_d\).  The inclusions among
\(q,qa,qc\) then prove all three equalities.
\end{proof}

For later equality arguments it is useful to retain the complete gap.  If
\(L_y=V_y\absop{L_y}\), define \(P_y=P_{L_y,B_y}\) and
\(G=(M_d\Id-F_d(U))^{1/2}\).  Then
\begin{equation}
\boxed{
 M_d\Id-\cI_d
 =\frac12\sum_yP_y^\dagger P_y
 +\frac12G^\dagger G
 +\frac12(GA_0^\dagger)^\dagger(GA_0^\dagger).
}
\label{eq:global-certificate}
\end{equation}
This is an operator identity, not a finite-level NPA certificate.
At the operator level, scalar functional calculus also gives
\[
 F_d(U)=M_d\Id
 \quad\Longleftrightarrow\quad
 U^d=(-1)^{d-1}\Id.
\]
For an individual maximizing vector \(\lvert\psi\rangle\), however, the
exact conclusion from \eqref{eq:global-certificate} is only
\[
 P_y\lvert\psi\rangle=0\ \ (0\leq y<d),\qquad
 G\lvert\psi\rangle=0,\qquad
 GA_0^\dagger\lvert\psi\rangle=0.
\]
We do not promote these support-level conditions to a global relation for
\(U\), or to a classification of all maximizing realizations.

\begin{corollary}[First augmented family]
\label{cor:first-augmented}
For
\[
 \overline{\cI}_d=\cI_d+\Reop(A_0B_d)
\]
with a further order-\(d\) Bob observable,
\[
 \beta_q(\overline{\cI}_d)
 =\beta_{qa}(\overline{\cI}_d)
 =\beta_{qc}(\overline{\cI}_d)=M_d+1.
\]
\end{corollary}

\begin{proof}
The upper bound follows from Theorem~\ref{thm:exact} and
\(\Reop(A_0B_d)\leq\Id\).  The finite strategy in
Appendix~\ref{app:attainment}, with the aligned added observable, attains
equality.
\end{proof}

\section{Equality phases and permutation blindness}
\label{sec:phase}

The scalar equality set in \eqref{eq:scalar-equality} is a set, not an
ordering.  A Bell score built only from first powers can therefore fail to
record how those phases are arranged along a weighted cycle.  The following
general statement isolates the exact sufficient conditions.  It does not
claim to classify every maximizer of every polar-linear functional.

Let \(R\) be a finite coefficient index set, let
\(g_r(z)=\alpha_r+\beta_rz\), and put
\[
 f(z)=\sum_{r\in R}\absop{g_r(z)},
 \qquad M=\max_{z\in\bbT}f(z).
\]
For order-\(d\) observables define
\begin{equation}
 \cJ=\sum_{r\in R}\Reop\!\left[(\alpha_rA_0+\beta_rA_1)B_r\right]
 +\Reop(A_0B_*).
 \label{eq:polar-linear-J}
\end{equation}

\begin{theorem}[Conditional phase-permutation theorem]
\label{thm:permutation}
Suppose there are labeled phases \(z_j,s_{rj}\in\bbT\),
\(0\leq j<d\), such that
\[
 f(z_j)=M,\qquad \prod_jz_j=1,
\]
and, for every \(r\),
\[
 s_{rj}\absop{g_r(z_j)}=g_r(z_j),
 \qquad \prod_js_{rj}=1.
\]
At a zero of \(g_r\), \(s_{rj}\) may be any unitary extension phase.
For every permutation \(\kappa\) of the labels, set on \(\bbC^d\)
\begin{equation}
\begin{aligned}
 A_0&=X,&
 A_1&=X\operatorname{diag}(z_{\kappa_j}),\\
 B_r&=\overline{X\operatorname{diag}(s_{r,\kappa_j})},&
 B_*&=X,
\end{aligned}
\label{eq:permutation-strategy}
\end{equation}
where \(X\lvert j\rangle=\lvert j+1\bmod d\rangle\), and use
\(\lvert\Phi_d\rangle=d^{-1/2}\sum_j\lvert j,j\rangle\).
Then:
\begin{enumerate}[label=(\roman*),leftmargin=1.7em]
\item every displayed observable is unitary, has \(d\)-th power \(\Id\),
and has the full simple \(d\)-th-root spectrum;
\item the strategy attains the operator upper bound \(M+1\) for
\eqref{eq:polar-linear-J}; and
\item all local first moments and every complex first-harmonic correlator
\(\langle A_xB_r\rangle\), \(\langle A_xB_*\rangle\) are independent of
\(\kappa\).
\end{enumerate}
\end{theorem}

\begin{proof}
With \(U=A_0^\dagger A_1\) and
\(L_r=A_0g_r(U)=V_r\absop{L_r}\), Lemma~\ref{lem:polar} and functional
calculus give
\[
 (M+1)\Id-\cJ
 =\frac12\sum_rP_r^\dagger P_r
 +\frac12R_*^\dagger R_*
 +\frac12\bigl(H+A_0HA_0^\dagger\bigr),
\]
where
\[
 P_r=\absop{L_r^\dagger}^{1/2}-V_r\absop{L_r}^{1/2}B_r,
 \quad R_*=\Id-A_0B_*,
 \quad H=M\Id-f(U)\geq0.
\]
Thus \(M+1\) is an operator upper bound.

For arbitrary nonzero weights \(t_j\),
\begin{equation}
 \left(X\operatorname{diag}(t_j)\right)^d
 =\left(\prod_jt_j\right)\Id,
 \label{eq:weighted-cycle}
\end{equation}
and its characteristic polynomial is
\(t^d-\prod_jt_j\).  The product hypotheses therefore prove (i).  Moreover,
\[
 \alpha_rA_0+\beta_rA_1
 =X\operatorname{diag}(g_r(z_{\kappa_j}))
\]
has unitary polar extension
\(X\operatorname{diag}(s_{r,\kappa_j})\).  The maximally-entangled trace
identity
\[
 \langle\Phi_d\vert R\otimes S\vert\Phi_d\rangle
 =\frac1d\Tr(RS^T)
\]
annihilates every factor in the displayed gap, proving (ii).  Finally,
\[
\begin{aligned}
 \langle A_0B_r\rangle&=\frac1d\sum_j\overline{s_{rj}},&
 \langle A_1B_r\rangle&=\frac1d\sum_jz_j\overline{s_{rj}},\\
 \langle A_0B_*\rangle&=1,&
 \langle A_1B_*\rangle&=\frac1d\sum_jz_j.
\end{aligned}
\]
These are symmetric sums of the labeled data.  Every local observable is a
weighted shift and has zero trace.  This proves (iii).
\end{proof}

The theorem explains both its power and its limit.  It gives a concrete
permutation orbit within a maximizing face, but it says nothing about
maximizers whose relative spectrum is different, has different
multiplicities, acts reducibly, or uses a kernel outside the stated phase
data.  Claims about the complete maximizing face require separate rigidity
analysis.

\section{Exact nonuniform maximizers of the first family}
\label{sec:biased}

Put
\[
 \delta_d=\begin{cases}0,&d\text{ odd},\\1,&d\text{ even},\end{cases}
 \qquad
 z_k=\exp\!\left(\frac{\pi i(2k+\delta_d)}d\right).
 \label{eq:equality-roots}
\]
These are precisely the roots in \eqref{eq:scalar-equality}.  For
\(0\leq y,k<d\), define
\[
 s_{y,k}=\frac{1+\omega^yz_k}{\absop{1+\omega^yz_k}}.
\]
The denominator never vanishes.  The polynomial
\(t^d-(-1)^{d-1}\) gives
\begin{equation}
 \prod_kz_k=1,
 \qquad
 \prod_k(1+\omega^yz_k)=2,
 \qquad
 \prod_ks_{y,k}=1.
 \label{eq:cyclic-products}
\end{equation}
Thus Theorem~\ref{thm:permutation} applies with
\(g_y(z)=1+\omega^yz\).

For every permutation \(\kappa\), the explicit first-family strategy is
\begin{equation}
\begin{aligned}
 A_0&=X,& A_1&=X Z_\kappa,&
 Z_\kappa&=\operatorname{diag}(z_{\kappa_j}),\\
 V_y&=X\operatorname{diag}(s_{y,\kappa_j}),&
 B_y&=\overline{V_y},& B_d&=X.
\end{aligned}
\label{eq:first-strategy}
\end{equation}
Each observable defines a rank-one projective \(d\)-outcome measurement and
\[
 \langle\overline{\cI}_d\rangle=M_d+1.
\]
The Bell-visible correlators are independent of the ordering.  Explicitly,
with \(\alpha_d=\csc(\pi/(2d))\),
\begin{equation}
 \langle A_0B_y\rangle=\frac{\alpha_d}{d},
 \qquad
 \langle A_1B_y\rangle=\frac{\alpha_d\omega^{-y}}d,
 \qquad
 \langle A_0B_d\rangle=1,
 \qquad
 \langle A_1B_d\rangle=0.
 \label{eq:visible-correlators}
\end{equation}
The cyclic ordering \(\kappa_j=j\) reproduces the canonical first-harmonic
behavior.  Other orderings can change higher powers and full probabilities.

To calculate the designated pair \((A_1,B_d)\), set
\begin{equation}
 q_0=1,
 \qquad q_{j+1}=z_{\kappa_j}q_j,
 \qquad
 \widehat q_m=\sum_{j=0}^{d-1}q_j\omega^{mj}.
 \label{eq:q-sequence}
\end{equation}
The product condition makes the sequence cyclic.  Direct diagonalization of
the two weighted shifts gives
\begin{equation}
 p_\kappa(a,b\mid1,d)
 =\frac{\absop{\widehat q_{-(a+b)}}^2}{d^3}.
 \label{eq:target-table}
\end{equation}
Parseval gives \(\sum_m\absop{\widehat q_m}^2=d^2\).  Consequently both
local marginals are exactly \(1/d\), while the joint table is uniform if and
only if every squared Fourier magnitude equals \(d\).

\begin{theorem}[Biased exact maximizers]
\label{thm:biased}
For every \(d\geq4\), choose the final-two swap
\begin{equation}
 \kappa=(0,1,\ldots,d-3,d-1,d-2).
 \label{eq:final-swap}
\end{equation}
The strategy \eqref{eq:first-strategy} attains \(M_d+1\), has uniform local
marginals, and has a nonuniform target table.  With one-dimensional Eve,
\begin{equation}
 G(AB\mid1,d,E)
 \geq\frac1{d^2}
 +\frac{2\sin(\pi/d)\sin(3\pi/d)}{d^2(d-1)}
 >\frac1{d^2}.
 \label{eq:guessing-gap}
\end{equation}
\end{theorem}

\begin{proof}
The second cyclic autocorrelation of \eqref{eq:q-sequence} is
\[
 R_2=\sum_jq_{j+2}\overline{q_j}
 =\sum_jz_{\kappa_j}z_{\kappa_{j+1}}.
\]
It vanishes for the cyclic ordering.  The swap \eqref{eq:final-swap} gives
\begin{equation}
 R_2=(z_{d-1}-z_{d-2})(z_{d-3}-z_0),
 \qquad
 \absop{R_2}=4\sin(\pi/d)\sin(3\pi/d)>0.
 \label{eq:R2}
\end{equation}
Hence the Fourier magnitudes are not all \(d\).

For the quantitative estimate put
\(x_m=\absop{\widehat q_m}^2-d\) and
\(\Delta=\max_mx_m\).  Fourier inversion at lag two gives
\[
 R_2=\frac1d\sum_mx_m\omega^{-2m},
\]
where the constant term vanishes for \(d\geq4\).  Since the nonzero
mean-zero sequence \((x_m)\) has at most \(d-1\) positive entries,
\[
 d\absop{R_2}\leq\sum_m\absop{x_m}
 =2\sum_{x_m>0}x_m\leq2(d-1)\Delta.
\]
With trivial Eve, always guessing a most likely pair gives
\(G=(d+\Delta)/d^3\).  Substituting \eqref{eq:R2} proves
\eqref{eq:guessing-gap}.
\end{proof}

For \(d=2,3\), every ordering in this particular permutation family is
Fourier-flat.  The theorem therefore neither proves nor disproves
Bell-value rigidity in those dimensions.

\section{Extension to the second augmented cyclic family}
\label{sec:second}

The second reduced family of Ref.~\cite{PeritoEtAl2026} has a different
Alice measurement set but compresses the same first-family Bob observables.
Put \(\eta=\e^{\pi i/d}\), so \(\omega=\eta^2\), and define
\begin{equation}
 \lambda_\ell=
 \frac{(-1)^{\ell-1}\eta^{\ell(\ell-1)}}
 {d\sin(\pi(\ell-\tfrac12)/d)},
 \qquad 0\leq\ell<d,
 \label{eq:lambda}
\end{equation}
where \((-1)^{-1}=-1\) at \(\ell=0\).  For Bob's observables in
\eqref{eq:first-strategy}, let
\[
 \widehat B_\ell=\sum_{y=0}^{d-1}\omega^{\ell y}B_y
\]
and define the Hermitian functional
\begin{equation}
 \cF_d=\sum_{\ell=0}^{d-1}
 \Reop\!\left(\overline{\lambda_\ell}A_\ell\widehat B_\ell\right),
 \qquad
 \overline{\cF}_d=\cF_d+\Reop(A_0B_d).
 \label{eq:second-functional}
\end{equation}

The source SOS is short and load-bearing.  Fourier orthogonality and
unitarity of the \(B_y\)'s give
\[
 \sum_\ell \widehat B_\ell^\dagger\widehat B_\ell=d^2\Id,
 \qquad \sum_\ell\absop{\lambda_\ell}^2=1.
\]
Thus, for
\(P_\ell=d\lambda_\ell\Id-A_\ell\widehat B_\ell\),
\begin{equation}
 d\Id-\cF_d=\frac1{2d}\sum_{\ell=0}^{d-1}P_\ell^\dagger P_\ell.
 \label{eq:second-sos}
\end{equation}
The identity requires only cross-party commutation, so the source value also
has the immediate commuting-operator reading
\[
 \beta_q(\cF_d)=\beta_{qa}(\cF_d)=\beta_{qc}(\cF_d)=d,
 \qquad
 \beta_q(\overline{\cF}_d)=\beta_{qa}(\overline{\cF}_d)
 =\beta_{qc}(\overline{\cF}_d)=d+1.
\]

\begin{theorem}[Permutation maximizers of the second family]
\label{thm:second}
For every permutation \(\kappa\), define
\begin{equation}
 r_\ell=\eta^{-\ell(\ell-1+\delta_d)},
 \qquad
 D_\ell=r_\ell X\operatorname{diag}
 (\omega^{-\ell\kappa_j}),
 \qquad
 A_\ell=\overline{D_\ell}.
 \label{eq:second-A}
\end{equation}
Together with the \(B_y\)'s in \eqref{eq:first-strategy}, these are
order-\(d\) observables and
\begin{equation}
 \widehat B_\ell=d\lambda_\ell D_\ell.
 \label{eq:Fourier-compression}
\end{equation}
They annihilate every factor in \eqref{eq:second-sos}.  With \(B_d=X\), the
augmented value is \(d+1\).  For the swap \eqref{eq:final-swap}, the target
pair \((A_1,B_d)\) has exactly the table \eqref{eq:target-table} and the gap
\eqref{eq:guessing-gap} for every \(d\geq4\).  The complete matrix of complex
first-harmonic correlators is independent of \(\kappa\).  Under the
alternative source-appendix convention that adjoints all Bob observables,
the target table is the same up to the harmless outcome inversion
\(b\mapsto-b\).
\end{theorem}

\begin{proof}
For \(z_t=\eta^{2t+\delta_d}\), write
\[
 s_{0,t}=\varepsilon_t\eta^{t+\delta_d/2},
 \qquad
 \varepsilon_t=\begin{cases}
 +1,&2t+\delta_d<d,\\
 -1,&2t+\delta_d>d.
 \end{cases}
\]
There is no equality case, and
\(s_{y,t}=s_{0,t+y\bmod d}\).  The diagonal weight of
\(\widehat B_\ell\) at a
label \(t\) is therefore \(\omega^{-\ell t}S_\ell\), with
\[
 S_\ell=\sum_t\omega^{\ell t}\overline{s_{0,t}}.
\]
Let \(n=(d+1-\delta_d)/2\) and \(u_\ell=\eta^{2\ell-1}\).  Since
\(u_\ell^d=-1\), two geometric sums give
\[
 S_\ell
 =\eta^{-\delta_d/2}
 \left(\sum_{t=0}^{n-1}u_\ell^t-\sum_{t=n}^{d-1}u_\ell^t\right)
 =d\lambda_\ell r_\ell.
\]
This is \eqref{eq:Fourier-compression}, including its phase.

By \eqref{eq:weighted-cycle} and
\(\sum_j\kappa_j=d(d-1)/2\),
\[
 D_\ell^d
 =(-1)^{\ell(\ell-1+\delta_d)+\ell(d-1)}\Id=\Id,
\]
because \(\ell(\ell+d-2+\delta_d)\) is even.  Equation
\eqref{eq:Fourier-compression} gives
\[
 P_\ell\lvert\Phi_d\rangle
 =d\lambda_\ell
 (\Id-\overline{D_\ell}\otimes D_\ell)\lvert\Phi_d\rangle=0.
\]
Also \(A_0=X\), so \(X\otimes X\) fixes \(\lvert\Phi_d\rangle\).
Finally
\[
 A_1=X\operatorname{diag}(\eta^{2\kappa_j+\delta_d})
 =X\operatorname{diag}(z_{\kappa_j}),
\]
which identifies the target table with \eqref{eq:target-table}.  Each
first-harmonic correlator is a sum over the complete unordered root set, and
is therefore permutation-independent.
\end{proof}

The theorem uses the source SOS to prove global optimality; merely
annihilating a candidate collection of operators would not have sufficed.
What is common to both families is the root ordering, the Bob cycle, and the
target table.  The second family additionally requires the exact Fourier
phases in \eqref{eq:second-A}.  No statement is made here about a complete
classification of its maximizing face.

\section{Randomness interpretation}
\label{sec:randomness}

Define the value-conditioned guessing quantity for either augmented family
by
\[
 G_{\mathrm{val}}(d)
 =\sup_{\mathsf S:\,\langle\mathcal B_d\rangle=\beta(\mathcal B_d)}
 G_{\mathsf S}(AB\mid1,d,E),
\]
where \(\mathcal B_d\) is \(\overline{\cI}_d\) or
\(\overline{\cF}_d\).  Theorems~\ref{thm:biased} and \ref{thm:second} imply
\begin{equation}
 G_{\mathrm{val}}(d)
 \geq\frac1{d^2}
 +\frac{2\sin(\pi/d)\sin(3\pi/d)}{d^2(d-1)}
 >\frac1{d^2}
 \qquad(d\geq4).
 \label{eq:value-conditioned}
\end{equation}
This is a lower bound on the worst case over the maximizing face.  We do not
claim that the final-two swap is worst possible.

The scope can be stated in one sentence:
\begin{quote}
\emph{The maximum scalar Bell value alone does not identify the canonical
behavior or certify maximal global randomness across the entire maximizing
face.}
\end{quote}
It is equally important to state what does not follow.  The result does not
dispute a calculation that fixes the complete canonical behavior, and it does
not show that the canonical strategy lacks maximal global randomness.  In
fact, the numerical program described in Appendix B.1 of
Ref.~\cite{PeritoEtAl2026} reconstructs the full distribution of the
canonical strategy before optimizing Eve's guessing probability.  The
root-swapped strategy has different higher Fourier correlators and therefore
does not satisfy those full-behavior constraints.

A full-behavior SDP may thus return \(G=1/d^2\) for the canonical table while
\eqref{eq:value-conditioned} holds for the scalar maximum.  Passing from the
first statement to the second would require uniqueness of the maximizing
behavior, self-testing strong enough to control the target PVMs, or other
statistics that exclude the phase-permuted realization.  The source's
\(d=3\) self-testing theorem for the second family is outside our
\(d\geq4\) construction and is not challenged.

\begin{corollary}[No value-only endpoint robustness]
For either augmented family and every \(d\geq4\), there is no valid bound for
all strategies of the form
\[
 G(AB\mid E)\leq\frac1{d^2}+f_d(\varepsilon),
 \qquad f_d(\varepsilon)\longrightarrow0,
\]
where \(\varepsilon\) is only the Bell deficit from the maximum.
More precisely, the asserted bound is required to hold for every strategy
whose deficit is at most \(\varepsilon\).
\end{corollary}

\begin{proof}
The biased exact maximizer has deficit zero, hence is feasible for every
\(\varepsilon\geq0\), and
has the fixed positive gap in \eqref{eq:value-conditioned}.  Bounds that use
additional observed correlators or the complete behavior are not excluded.
\end{proof}

\section{Setting-complexity consequences}
\label{sec:settings}

The permutation mechanism is a design warning, not a universal lower bound
on settings.  One universal baseline is nevertheless immediate.

\begin{proposition}[One-input baseline]
\label{prop:one-input}
If either party has only one input, every finite-output nonsignalling behavior
is local and admits a compatible finite-dimensional pure projective
realization in which Eve guesses both outputs perfectly.  Hence at least two
inputs per party are necessary for any private global randomness, even when
the complete behavior is fixed.
\end{proposition}

\begin{proof}
Suppose Alice has one input.  Nonsignalling makes
\(p(a)=\sum_bp(a,b\mid y)\) independent of \(y\).  For \(p(a)>0\), let
\(r_y(b\mid a)=p(a,b\mid y)/p(a)\), and choose an arbitrary conditional
distribution when \(p(a)=0\).  The hidden variable
\[
 \lambda=(a,b_1,\ldots,b_{m_B}),
 \qquad
 \mu(\lambda)=p(a)\prod_yr_y(b_y\mid a)
\]
reproduces the whole behavior by outputting the stored entries.  The pure
state
\[
 \sum_\lambda\sqrt{\mu(\lambda)}
 \lvert\lambda\rangle_A\lvert\lambda\rangle_B
 \lvert\lambda\rangle_E
\]
with diagonal grouping PVMs realizes it.  Eve reads \(\lambda\) and guesses
perfectly.  The other case is symmetric.
\end{proof}

This is a standard locality fact rather than a novelty claim.  For \(d=2\),
the \((2,2)\) baseline is attainable by a known tilted Bell score
\cite{WooltortonEtAl2022}.  For \(d\geq3\), neither this paper nor the
phase-permutation theorem determines the minimum number of settings.  A
current all-dimensional construction with projective generation
measurements uses two Alice settings and \(d^2+1\) Bob settings
\cite{FarkasMironowiczAugusiak2026}.

Appendix~\ref{app:settings} records two narrower observations retained from
the setting-complexity study: the standard two-input qudit self-test has no
uniform joint setting pair, and one natural separately bounded third-setting
repair cannot expose its common computational mutually unbiased basis.  The
obstruction does not rule out a genuinely joint SOS, another MUB, another
self-test, or an arbitrary \((2,3,d,d)\) Bell functional.

\section{Discussion and open problems}

The exact value theorem and the biased maximizers expose a structural
distinction between optimality and rigidity.  The first operator controls a
sum of moduli of first-order linear functions of the relative unitary.  Its
score is sharp enough to determine the allowable scalar phases but not their
cyclic order.  That missing order is invisible to the Bell score and visible
to the target output distribution.

Several questions are now sharply posed.
\begin{enumerate}[leftmargin=1.6em]
\item Characterize the complete maximizing faces of both augmented families,
including reducible and commuting-operator realizations.
\item Determine the exact worst-case guessing probability at the Bell maximum;
the bound \eqref{eq:value-conditioned} need not be tight.
\item Resolve the \(d=2,3\) maximizing faces for the first augmented family.
The permutation orbit studied here is uniform in those dimensions.
\item Classify which phase permutations are related by local isometries and
output relabelings, and which produce genuinely different full behaviors.
\item Find a low-setting modification whose Bell score controls the missing
higher Fourier data, and determine whether it is rigid or robustly
self-testing.
\item Develop robust bounds conditioned on a minimal set of extra correlators,
rather than either one scalar score or the entire behavior.
\end{enumerate}
These problems are natural opportunities for collaboration with the
originating authors and the broader Bell-inequality community.  No external
review or collaboration is claimed here.

\section*{Data, code, and historical versions}

The accompanying repository contains the TeX source, exact and numerical
verifiers, machine-readable certificates, claims ledger, proof-dependency
map, fresh adversarial and priority audits, and integrity hashes.  One command
replays the mathematical checks, the two independent \(d=4\) exact
implementations, the manuscript build, and the website validation.  The
three earlier standalone papers and their original PDFs, sources, hashes,
and Git history remain available as historical artifacts; they are not
silently replaced by this synthesis.

\section*{AI-assistance disclosure}

Generative-AI systems assisted with proof exploration, adversarial testing,
code, literature-search organization, and drafting under Alec Kriebel's
direction.  Alec Kriebel is the sole named author and takes responsibility
for the manuscript and release.  AI output was not treated as evidence:
claims were retained only when accompanied by a written proof and replayable
checks.  No originating author or outside researcher was contacted.

\appendix

\section{Proof of the scalar extremum}
\label{app:scalar}

Write \(z=\e^{2is}\), \(h=\pi/d\), and \(a=h/2\).  Half of the sum in
\eqref{eq:scalar-bound} is
\[
 G_d(s)=\sum_{y=0}^{d-1}\absop{\cos(s+yh)},
\]
which has period \(h\).

If \(d=2m+1\), take \(-a\leq s\leq a\) and reindex by
\(k=-m,\ldots,m\).  Every angle lies in \([-\pi/2,\pi/2]\), so
\[
 G_d(s)=\sum_{k=-m}^m\cos(s+kh)=\csc(a)\cos s.
\]
The maximum occurs exactly at \(s\equiv0\pmod h\), equivalently \(z^d=1\).

If \(d=2m\), take \(0\leq s\leq h\).  The first \(m\) cosines are
nonnegative and the last \(m\) are nonpositive.  Two finite sums give
\[
 G_d(s)=\sum_{y=0}^{m-1}\cos(s+yh)
 -\sum_{y=m}^{2m-1}\cos(s+yh)
 =\csc(a)\cos(s-a).
\]
The maximum occurs exactly at \(s\equiv a\pmod h\), equivalently \(z^d=-1\).
Multiplying by two proves Lemma~\ref{lem:scalar}, including the equality set.

\section{Rechecking the canonical attaining strategy}
\label{app:attainment}

This appendix verifies the lower strategy supplied in
Ref.~\cite{PeritoEtAl2026} in a polar form.  On \(\bbC^d\), let
\[
 Z\lvert j\rangle=\omega^j\lvert j\rangle,
 \qquad X\lvert j\rangle=\lvert j+1\bmod d\rangle,
 \qquad A_0=Z,\quad A_1=X.
\]
Set
\[
 W_y=\omega^yZ^\dagger X,
 \qquad H_y=\absop{\Id+W_y},
 \qquad
 Q_y=H_y^{-1}(\Id+W_y^\dagger)Z^\dagger,
 \qquad B_y=Q_y^T.
\]
The weighted-shift action gives
\[
 W_y^d=(-1)^{d-1}\Id,
 \qquad
 \spec(W_y)=\{\lambda:\lambda^d=(-1)^{d-1}\}.
\]
The spectrum is simple, and \(-1\) is absent; hence \(H_y\) is invertible.
For \(L_y=Z+\omega^yX=Z(\Id+W_y)\),
\[
 L_y=\underbrace{Z(\Id+W_y)H_y^{-1}}_{V_y}H_y
\]
is the polar decomposition, and \(Q_y=V_y^\dagger\) is unitary.

It remains to check the measurement order.  On the displayed spectrum put
\(p(z)=(1+z)/\absop{1+z}\).  Since
\(Z^\dagger W_yZ=\omega^{-1}W_y\),
\[
 Q_y^d=\prod_{r=0}^{d-1}p(\omega^{-r}W_y)^\dagger.
\]
For each eigenvalue \(\lambda\), the corresponding scalar equals
\[
 \overline{\frac{\prod_r(1+\omega^{-r}\lambda)}
 {\prod_r\absop{1+\omega^{-r}\lambda}}}
 =\overline{\frac{1-(-\lambda)^d}{\absop{1-(-\lambda)^d}}}=1.
\]
Thus \(B_y^d=\Id\).  The weighted-cycle characteristic polynomial shows
that each \(d\)-th root occurs.

Finally,
\[
 \langle\Phi_d\vert L_y\otimes B_y\vert\Phi_d\rangle
 =\frac1d\Tr(L_yV_y^\dagger)=\frac1d\Tr(H_y).
\]
The spectrum of \(W_y\) is the scalar equality set, so summing over \(y\)
gives \(\langle\cI_d\rangle=M_d\).  Taking the extra Bob observable
\(B_d=Z^\dagger\) gives \(\langle Z\otimes Z^\dagger\rangle=1\), proving
Corollary~\ref{cor:first-augmented}.  This appendix claims no new attaining
strategy.

\section{The exact four-outcome certificate}
\label{app:d4}

Let \(d=4\), \(\zeta=\e^{\pi i/8}\), and
\(\kappa=(0,1,3,2)\).  Then
\[
 A_0=X,
 \qquad
 A_1=X\operatorname{diag}(\zeta^2,\zeta^6,\zeta^{14},\zeta^{10}),
 \qquad B_4=X.
\]
The first-family polar unitaries
\(V_y=X\operatorname{diag}(\zeta^{e_{y,0}},\ldots,\zeta^{e_{y,3}})\)
are encoded by
\[
\begin{array}{c@{\qquad}cccc}
\toprule
y&e_{y,0}&e_{y,1}&e_{y,2}&e_{y,3}\\
\midrule
0&1&3&15&13\\
1&3&13&1&15\\
2&13&15&3&1\\
3&15&1&13&3\\
\bottomrule
\end{array}
\qquad B_y=\overline{V_y}.
\]
These are exact monomial matrices over
\(\mathbb Q(\zeta_{16})\), not numerical polar decompositions.

The recurrence \eqref{eq:q-sequence} gives
\[
 q=(1,\zeta^2,-1,\zeta^6),
 \qquad
 (\absop{\widehat q_m}^2)_{m=0}^3=(2,6,2,6).
\]
For both augmented families,
\begin{equation}
 p(a,b\mid1,4)=
 \begin{cases}
 1/32,&a+b\text{ even},\\
 3/32,&a+b\text{ odd},
 \end{cases}
 \qquad G=3/32>1/16.
 \label{eq:d4-table}
\end{equation}
One exact verifier reconstructs \eqref{eq:d4-table} from projectors and a
separate Fourier calculation.  An independent exact implementation
reconstructs the source coefficients \(\lambda_{y,k}\), proves
\(\widehat B_\ell=4\lambda_\ell D_\ell\), checks all fourth-power relations,
expands
the full second-family SOS, and verifies values \(M_4+1\) and \(5\).

\section{Scoped low-setting observations}
\label{app:settings}

The standard two-input maximally-entangled-qudit self-test
\cite{SarkarEtAl2021} uses Fourier-phase bases.  Up to input swaps,
conjugation, transposition, and outcome relabeling, its four ideal tables have
the form
\begin{equation}
 p(a,b\mid x,y)
 =\frac{\sin^2(\pi\delta_{xy})}
 {d^3\sin^2(\pi(a-b+\delta_{xy})/d)},
 \qquad
 (\delta_{xy})=
 \begin{pmatrix}1/4&3/4\\-1/4&1/4\end{pmatrix}.
 \label{eq:standard-tables}
\end{equation}
Every table has
\[
 p_{\max}=\frac1{2d^3\sin^2(\pi/(4d))}>\frac1{d^2},
\]
so no ideal input pair is uniform.  Perfectly anchoring a third Bob setting
to either Alice basis leaves a cross table with
\[
 p_{\max}^{\mathrm{cross}}
 =\frac1{d^3\sin^2(\pi/(2d))}>\frac1{d^2}
 \qquad(d\geq3).
\]
The exception \(d=2\) is precisely the mutually-unbiased qubit case.

There is a broader but still scoped obstruction.  Let \(F\) be the Fourier
matrix, \(\rho=\e^{\pi i/d}\), and
\(R=\operatorname{diag}(1,\rho,\ldots,\rho^{d-1})\).  The real Hermitian
operator system generated by the two ideal Alice PVMs is
\[
 \mathcal S=\{F\operatorname{diag}(x)F^\dagger
 +RF\operatorname{diag}(y)F^\dagger R^\dagger:x,y\in\mathbb R^d\}.
\]

\begin{proposition}[Computational-MUB exposure obstruction]
\label{prop:mub}
If \(K\in\mathcal S\) has a computational vector \(e_r\) as an eigenvector
with eigenvalue \(\kappa\), then either \(K=\kappa\Id\), or \(\kappa\) is
neither the largest nor the smallest eigenvalue of \(K\).  Consequently, no
nontrivial separately bounded term
\[
 \mathcal L=\sum_bK_b\otimes Q_b,
 \qquad K_b\in\mathcal S,
\]
can expose the computational PVM on \(\lvert\Phi_d\rangle\) through a
coefficientwise spectral upper bound.
\end{proposition}

\begin{proof}
Write \(K=C+RC'R^\dagger\) with \(C,C'\) Hermitian circulants.  Comparing
the column-\(r\) eigenvector equations with their Hermitian conjugates forces
all but the first and last
\(s=\min(r,d-1-r)\) cyclic diagonals to vanish.  After grouping the first
\(s\), middle \(d-2s\), and last \(s\) coordinates,
\[
 K-\kappa\Id=
 \begin{pmatrix}0&0&T\\0&0&0\\T^\dagger&0&0\end{pmatrix},
\]
where \(T\) is upper-triangular Toeplitz.  Its nonzero eigenvalues occur in
pairs \(\pm\sigma_j(T)\).  Hence a nonscalar \(K\) has spectrum on both
sides of \(\kappa\).

To saturate a coefficientwise upper bound on the maximally entangled state
with computational \(Q_b\), every computational outcome has nonzero weight
and the corresponding \(e_b\) must be a maximal eigenvector of \(K_b\).
The first part makes that impossible unless all relevant \(K_b\) are scalar,
in which case the term constrains no measurement.
\end{proof}

Proposition~\ref{prop:mub} rules out only this separately bounded local-
spectral route.  It is not a no-go theorem for \((2,3,d,d)\), and it is not
used in any cyclic-family conclusion.

\section{Source conventions and claim boundary}
\label{app:conventions}

Our first-family definition \eqref{eq:Id} follows the displayed main-text
operator definition in Ref.~\cite{PeritoEtAl2026}.  A nearby appendix writes
one Bob adjoint differently while its subsequent calculation follows the
main convention.  Replacing every Bob observable by its adjoint transports
between the two notations, but mixing them within one formula does not.  We
therefore state every adjoint explicitly and test the Fourier orientation.

For the second family, \eqref{eq:second-functional} matches the Hermitian
content of the source definition and its corrected \(1/(2d)\) SOS factor in
the latest version.  The source text adjacent to the canonical construction
briefly omits Alice's entrywise conjugation, while its theorem includes it.
Our choice \(A_\ell=\overline{D_\ell}\) is the one that annihilates the
displayed SOS on \(\lvert\Phi_d\rangle\).

If instead every Bob observable in the source appendix is written with an
adjoint, set \(B'_y=B_y^\dagger\) and \(B'_d=B_d^\dagger\).  This is Bob's
outcome inversion: it sends \(p(a,b)\) to \(p(a,-b)\) while preserving the
maximum, nonuniformity, and guessing probability.  It cannot be mixed
termwise with the no-adjoint convention used in
\eqref{eq:second-functional}.

The source's numerical guessing-probability table explicitly fixes the full
canonical distribution.  Our scalar-value conclusion should therefore not
be read as a criticism of that calculation.  The literal maximal-violation
conjecture and the fixed-behavior numerical experiment have different logical
hypotheses; the present construction separates them.

\begin{thebibliography}{99}
\raggedright

\bibitem{PeritoEtAl2026}
I. Perito, R. D'Avino, M. Jung, P. Mironowicz, A. Ac\'in, and R. Augusiak,
\newblock Bell inequalities tailored to optimal global randomness
certification,
\newblock \emph{arXiv:2606.21362v3} (2026),
\newblock \url{https://arxiv.org/abs/2606.21362v3}.

\bibitem{DharaEtAl2013}
C. Dhara, G. Prettico, and A. Ac\'in,
\newblock Maximal quantum randomness in Bell tests,
\newblock \emph{Phys. Rev. A} \textbf{88}, 052116 (2013),
\newblock \url{https://doi.org/10.1103/PhysRevA.88.052116}.

\bibitem{NietoSillerasEtAl2014}
O. Nieto-Silleras, S. Pironio, and J. Silman,
\newblock Using complete measurement statistics for optimal
device-independent randomness evaluation,
\newblock \emph{New J. Phys.} \textbf{16}, 013035 (2014),
\newblock \url{https://doi.org/10.1088/1367-2630/16/1/013035}.

\bibitem{BancalEtAl2014}
J.-D. Bancal, L. Sheridan, and V. Scarani,
\newblock More randomness from the same data,
\newblock \emph{New J. Phys.} \textbf{16}, 033011 (2014),
\newblock \url{https://doi.org/10.1088/1367-2630/16/3/033011}.

\bibitem{NPA2007}
M. Navascu\'es, S. Pironio, and A. Ac\'in,
\newblock Bounding the set of quantum correlations,
\newblock \emph{Phys. Rev. Lett.} \textbf{98}, 010401 (2007),
\newblock \url{https://doi.org/10.1103/PhysRevLett.98.010401}.

\bibitem{BuhrmanMassar2005}
H. Buhrman and S. Massar,
\newblock Causality and Tsirelson's bounds,
\newblock \emph{Phys. Rev. A} \textbf{72}, 052103 (2005),
\newblock \url{https://doi.org/10.1103/PhysRevA.72.052103}.

\bibitem{SalavrakosEtAl2017}
A. Salavrakos, R. Augusiak, J. Tura, P. Wittek, A. Ac\'in, and S. Pironio,
\newblock Bell inequalities tailored to maximally entangled states,
\newblock \emph{Phys. Rev. Lett.} \textbf{119}, 040402 (2017),
\newblock \url{https://doi.org/10.1103/PhysRevLett.119.040402}.

\bibitem{KaniewskiEtAl2019}
J. Kaniewski, I. \v{S}upi\'c, J. Tura, F. Baccari, A. Salavrakos, and
R. Augusiak,
\newblock Maximal nonlocality from maximal entanglement and mutually unbiased
bases, and self-testing of two-qutrit quantum systems,
\newblock \emph{Quantum} \textbf{3}, 198 (2019),
\newblock \url{https://doi.org/10.22331/q-2019-10-24-198}.

\bibitem{SarkarEtAl2021}
S. Sarkar, D. Saha, J. Kaniewski, and R. Augusiak,
\newblock Self-testing quantum systems of arbitrary local dimension with
minimal number of measurements,
\newblock \emph{npj Quantum Information} \textbf{7}, 151 (2021),
\newblock \url{https://doi.org/10.1038/s41534-021-00490-3}.

\bibitem{WooltortonEtAl2022}
L. Wooltorton, P. Brown, and R. Colbeck,
\newblock Tight analytic bound on the trade-off between device-independent
randomness and nonlocality,
\newblock \emph{Phys. Rev. Lett.} \textbf{129}, 150403 (2022),
\newblock \url{https://doi.org/10.1103/PhysRevLett.129.150403}.

\bibitem{BarizienEtAl2024}
V. Barizien, P. Sekatski, and J.-D. Bancal,
\newblock Custom Bell inequalities from formal sums of squares,
\newblock \emph{Quantum} \textbf{8}, 1333 (2024),
\newblock \url{https://doi.org/10.22331/q-2024-05-02-1333}.

\bibitem{KlepEtAl2026}
I. Klep, N. Leijenhorst, and V. Magron,
\newblock Robust self-testing with CHSH mod 3,
\newblock \emph{arXiv:2604.03700v1} (2026),
\newblock \url{https://arxiv.org/abs/2604.03700v1}.

\bibitem{CocciaEtAl2026}
L. Coccia, M. Padovan, and G. Vallone,
\newblock Systematic derivation of Tsirelson bounds in arbitrary dimensions,
\newblock \emph{arXiv:2606.21626v1} (2026),
\newblock \url{https://arxiv.org/abs/2606.21626v1}.

\bibitem{FarkasMironowiczAugusiak2026}
M. Farkas, P. Mironowicz, and R. Augusiak,
\newblock Maximal global device-independent randomness from projective
measurements in every dimension,
\newblock \emph{arXiv:2606.21369v2} (2026),
\newblock \url{https://arxiv.org/abs/2606.21369v2}.

\bibitem{DavinoEtAl2026}
R. D'Avino, I. Perito, P. Mironowicz, A. Ac\'in, and R. Augusiak,
\newblock Noise robustness of three outcome Bell certified quantum randomness,
\newblock \emph{arXiv:2606.21371v2} (2026),
\newblock \url{https://arxiv.org/abs/2606.21371v2}.

\end{thebibliography}

\end{document}
